% ============================================================================
% MOTEUR — COURS 07 : ÉLECTRONIQUE
% ============================================================================
\RequirePackage{silence}
\WarningFilter{xcolor}{Package option `usenames' is obsolete}

\documentclass[
  fontsize=10pt,
  twoside=true,
  open=any,
]{kaobook}

\usepackage[french]{babel}

\makeatletter
\def\input@path{{../../Style/}}
\makeatother
\NeedsTeXFormat{LaTeX2e}
\ProvidesPackage{TLCours}[2026/08/01 Socle commun des polycopiés de cours TL]

% -----------------------------------------------------------------------------
% Dépendances réellement communes aux cours
% -----------------------------------------------------------------------------
\RequirePackage{etoolbox}
\RequirePackage{needspace}
\RequirePackage[normalem]{ulem}

% Macros scientifiques et environnements historiques du framework.
\input{macros_SII_v2}
\input{environment_v2}

% Commandes personnalisées TL et couverture des cours.
\RequirePackage{CouvertureCours}
\RequirePackage{TL}

% Version explicite de PGFPlots, chargé par environment_v2.
\pgfplotsset{compat=1.18}

% -----------------------------------------------------------------------------
% Valeurs communes et compatibilité avec les anciens fichiers
% -----------------------------------------------------------------------------
\providecommand{\discipline}{Sciences Industrielles de l'Ingénieur}
\providecommand{\auteur}{Thomas Lusseau}
\providecommand{\repRel}{../../}
\providecommand{\repStyle}{\repRel/Style}
\providecommand{\ul}[1]{\uline{#1}}

\ifdefined\frenchsetup
  \frenchsetup{StandardItemLabels=true}
\fi

\graphicspath{{images/}{./}}

% Éviter qu'un titre soit isolé en bas de page.
\pretocmd{\section}{\Needspace{6\baselineskip}}{}{}
\pretocmd{\subsection}{\Needspace{5\baselineskip}}{}{}
\pretocmd{\subsubsection}{\Needspace{4\baselineskip}}{}{}

% Les hooks anglais de varioref ne sont pas utilisés dans les cours français.
\makeatletter
\@ifundefined{extrasenglish}{}{\let\extrasenglish\@empty}
\@ifundefined{noextrasenglish}{}{\let\noextrasenglish\@empty}
\makeatother

% ============================================================================
% BIBLIOTHÈQUE TL — ÉLECTRONIQUE ET ACQUISITION DU SIGNAL
% ============================================================================
\NeedsTeXFormat{LaTeX2e}
\ProvidesPackage{TLElectronique}[2026/08/03 v2026.08.03c Electronique et acquisition]
\RequirePackage[siunitx,european]{circuitikz}
\RequirePackage{xparse}
\RequirePackage{tikz}
\usetikzlibrary{positioning,arrows.meta,calc}
\typeout{TL-ELECTRONIQUE-MARKER-20260803-C}

\newcommand{\TLChaineAcquisition}{%
\begin{center}
\begin{tikzpicture}[node distance=7mm,>=Latex,font=\small]
\tikzset{bloc/.style={draw,rounded corners=1mm,minimum width=25mm,minimum height=10mm,align=center,fill=blue!4}}
\node[bloc] (cap) {Capteur};
\node[bloc,right=of cap] (cond) {Conditionnement};
\node[bloc,right=of cond] (ech) {Échantillonnage};
\node[bloc,right=of ech] (can) {Conversion\\A/N};
\node[bloc,right=of can] (num) {Traitement\\numérique};
\draw[->] (cap)--(cond); \draw[->] (cond)--(ech); \draw[->] (ech)--(can); \draw[->] (can)--(num);
\end{tikzpicture}
\end{center}}

% Version volontairement élémentaire du pont diviseur.
% circuitikz reste chargé et utilisé dans les autres schémas de la bibliothèque.
\newcommand{\TLPontDiviseur}{%
\typeout{TL-PONT-DIVISEUR-VERSION-C}%
\begin{center}
\begin{tikzpicture}[line width=0.8pt,font=\small]
  % Source représentée par une pile, sans calcul de coordonnées.
  \draw (0,0) -- (0,1.4);
  \draw (-0.45,1.4) -- (0.45,1.4);
  \draw (-0.25,1.7) -- (0.25,1.7);
  \draw (0,1.7) -- (0,4);
  \node[left] at (0,2.55) {$E$};

  % Conducteur supérieur.
  \draw (0,4) -- (2.5,4);

  % Résistance R1, symbole rectangulaire IEC.
  \draw (2.5,4) -- (2.5,3.6);
  \draw (2.15,3.6) rectangle (2.85,2.6);
  \node[right] at (2.85,3.1) {$R_1$};
  \draw (2.5,2.6) -- (2.5,2);

  % Sortie.
  \fill (2.5,2) circle (1.3pt);
  \draw (2.5,2) -- (4.2,2);
  \draw (4.2,2) circle (2.2pt);
  \node[right] at (4.2,2) {$u_s$};

  % Résistance R2.
  \draw (2.5,2) -- (2.5,1.6);
  \draw (2.15,1.6) rectangle (2.85,0.6);
  \node[right] at (2.85,1.1) {$R_2$};
  \draw (2.5,0.6) -- (2.5,0);
  \draw (2.5,0) -- (0,0);
\end{tikzpicture}
\end{center}}

\newcommand{\TLFiltreRCPasseBas}{%
\begin{center}
\begin{circuitikz}[european resistors]
\node[ocirc] at (0,2) {};
\node[left] at (0,2) {$u_e$};
\draw (0,2) -- (1,2) to[R,l=$R$] (4,2) -- (5,2);
\node[ocirc] at (5,2) {};
\node[right] at (5,2) {$u_s$};
\draw (4,2) to[C,l=$C$] (4,0) node[ground]{};
\draw (1,0) node[ground]{} -- (0,0);
\node[ocirc] at (0,0) {};
\end{circuitikz}
\end{center}}

\newcommand{\TLFiltreRCPasseHaut}{%
\begin{center}
\begin{circuitikz}[european resistors]
\node[ocirc] at (0,2) {};
\node[left] at (0,2) {$u_e$};
\draw (0,2) -- (1,2) to[C,l=$C$] (4,2) -- (5,2);
\node[ocirc] at (5,2) {};
\node[right] at (5,2) {$u_s$};
\draw (4,2) to[R,l=$R$] (4,0) node[ground]{};
\end{circuitikz}
\end{center}}

\newcommand{\TLAOPInverseur}{%
\begin{center}
\begin{circuitikz}[american voltages]
\draw (4,2) node[op amp] (op) {};
\node[ocirc] at (0,1.5) {};
\node[left] at (0,1.5) {$u_e$};
\draw (0,1.5) -- (1,1.5) to[R,l=$R_1$] (op.-);
\draw (op.+) -- ++(0,-1) node[ground]{};
\draw (op.out) -- ++(1.5,0) coordinate (sortieinv);
\node[ocirc] at (sortieinv) {};
\node[right] at (sortieinv) {$u_s$};
\draw (op.out) -- ++(0,1.5) to[R,l=$R_2$] ++(-3.5,0) |- (op.-);
\end{circuitikz}
\end{center}}

\newcommand{\TLAOPNonInverseur}{%
\begin{center}
\begin{circuitikz}[american voltages]
\draw (4,2) node[op amp] (op) {};
\node[ocirc] at (0,2.5) {};
\node[left] at (0,2.5) {$u_e$};
\draw (0,2.5) -- (op.+);
\draw (op.out) -- ++(1.5,0) coordinate (sortienoninv);
\node[ocirc] at (sortienoninv) {};
\node[right] at (sortienoninv) {$u_s$};
\draw (op.-) -- ++(-1,0) coordinate (n);
\draw (n) to[R,l_=$R_1$] ++(0,-2) node[ground]{};
\draw (op.out) -- ++(0,1.5) to[R,l=$R_2$] ++(-3.5,0) |- (n);
\end{circuitikz}
\end{center}}

\newcommand{\TLPontWheatstone}{%
\begin{center}
\begin{circuitikz}[european resistors]
\draw (0,0) -- (0,1.4);
\draw (-0.45,1.4) -- (0.45,1.4);
\draw (-0.25,1.7) -- (0.25,1.7);
\draw (0,1.7) -- (0,4) -- (2,4);
\node[left] at (0,2.55) {$E$};
\draw (2,4) to[R,l=$R_1$] (2,2) to[R,l=$R_2$] (2,0) -- (0,0);
\draw (2,4) -- (5,4) to[R,l=$R_3$] (5,2) to[R,l=$R_4$] (5,0) -- (2,0);
\draw (2,2) -- (3,2);
\node[ocirc] at (3,2) {};
\node[below] at (3,2) {$A$};
\draw (5,2) -- (4,2);
\node[ocirc] at (4,2) {};
\node[below] at (4,2) {$B$};
\draw[<->] (3,2.35)--(4,2.35) node[midway,above]{$u_{AB}$};
\end{circuitikz}
\end{center}}

\newcommand{\TLEchantillonnageSignal}{%
\begin{center}
\begin{tikzpicture}[x=.7cm,y=.8cm,>=Latex]
\draw[->] (0,0)--(10,0) node[right]{$t$};
\draw[->] (0,-1.4)--(0,2) node[above]{$u$};
\draw[domain=0:9.5,samples=200,smooth] plot(\x,{sin(50*\x)});
\foreach \x in {0,0.8,...,9.6}{\draw (\x,0)--(\x,{sin(50*\x)}) node[circle,fill,inner sep=1.2pt]{};}
\end{tikzpicture}
\end{center}}

\endinput

\input{Notations_Elec/Notations_Elec.sty}

% Environnement utilisé dans les sources du cours.
\newenvironment{application}
  {\par\medskip\noindent\textbf{Application.}\enspace}
  {\par\medskip}

% Métadonnées communes du polycopié.
\renewcommand{\discipline}{Sciences Industrielles de l'Ingénieur}
\renewcommand{\auteur}{Thomas Lusseau}

\raggedbottom
\setlength{\parindent}{0pt}
\setlength{\emergencystretch}{3em}
\makeindex[columns=3,title=Index alphabétique,intoc]

\begin{document}

% -----------------------------------------------------------------------------
% Couverture dans la charte des autres polycopiés
% -----------------------------------------------------------------------------
\begin{titlepage}
  \thispagestyle{empty}
  \begin{tikzpicture}[remember picture,overlay]
    \fill[TLCoverBlue]
      (current page.north west)
      rectangle ([yshift=-68mm]current page.north east);

    \fill[TLCoverFooterBackground]
      (current page.south west)
      rectangle ([yshift=20mm]current page.south east);

    \node[
      anchor=north west,
      xshift=15mm,
      yshift=-17mm,
      text=white,
      align=left,
      text width=175mm,
      font=\sffamily\bfseries\fontsize{25}{30}\selectfont
    ] at (current page.north west)
      {Acquisition et traitement\\du signal};

    \node[
      anchor=north west,
      xshift=16mm,
      yshift=-75mm,
      text=TLCoverOrange,
      font=\sffamily\bfseries\fontsize{15}{18}\selectfont
    ] at (current page.north west)
      {Cours 07 — Électronique};

    \node[
      anchor=north west,
      xshift=16mm,
      yshift=-94mm,
      text=TLCoverBlue,
      align=left,
      text width=170mm,
      font=\sffamily\fontsize{12}{17}\selectfont
    ] at (current page.north west)
      {Chaîne d'acquisition — Capteurs — Électrocinétique\\
       Représentation spectrale — Filtrage — Numérisation};

    \draw[TLCoverGreen,line width=2.5pt]
      ([xshift=16mm,yshift=-128mm]current page.north west)
      --
      ([xshift=-16mm,yshift=-128mm]current page.north east);

    \node[
      anchor=north west,
      xshift=16mm,
      yshift=-142mm,
      text=TLCoverBlue,
      align=left,
      text width=170mm,
      font=\sffamily\bfseries\fontsize{15}{19}\selectfont
    ] at (current page.north west)
      {ATS — Sciences Industrielles de l'Ingénieur};

    \node[
      anchor=south east,
      xshift=-12mm,
      yshift=8mm,
      text=TLCoverBlue,
      align=right,
      font=\sffamily\bfseries\fontsize{11}{13}\selectfont
    ] at (current page.south east)
      {Lycée Robert Doisneau — Thomas Lusseau};
  \end{tikzpicture}
  \null
\end{titlepage}

% Pied de page commun à tous les polycopiés.
\TLCourseFooter

\frontmatter
\tableofcontents
\clearpage

\mainmatter
\setchapterstyle{kao}
\setcounter{margintocdepth}{\sectiontocdepth}
\marginlayout
\pagestyle{scrheadings}

\chapter{Acquisition et traitement du signal}
\label{chap:electronique}
\margintoc

% Parties 1 à 4 : chaîne d'acquisition, capteurs et électrocinétique.
% ============================================================================
% COURS 07 — ACQUISITION ET TRAITEMENT DE L'INFORMATION
% Reprise fidèle du document 11-AcquisitionTraitementSignal2022.docx
% ============================================================================

\section*{Objectifs du cours}

Ce cours donne les démarches et les outils nécessaires à l'étude des chaînes d'information.

\subsection*{Savoirs}

À l'issue du cours, il faut connaître :
\begin{itemize}
  \item la structure d'une chaîne d'acquisition : capteur, conditionneur, amplificateur, filtre, convertisseur analogique-numérique, unité de traitement et convertisseur numérique-analogique ;
  \item la différence entre résolution et précision ;
  \item les impédances et admittances des dipôles de base ;
  \item les différents types de filtres ;
  \item la méthode permettant de déterminer rapidement le comportement d'un filtre aux basses et aux hautes fréquences ;
  \item le théorème de Shannon-Nyquist ;
  \item le filtre anti-repliement et son utilité ;
  \item la définition du quantum d'un CAN et d'un CNA ;
  \item les méthodes permettant de déterminer la résolution d'un CAN ou d'un CNA ;
  \item les bases 2, 10 et 16 et les méthodes de changement de base ;
  \item le codage en complément à deux et ses particularités ;
  \item les opérations simples en base 2 ;
  \item la notion de temps de conversion d'un CAN et sa relation avec la période d'échantillonnage.
\end{itemize}

\subsection*{Savoir-faire}

Il faut savoir :
\begin{itemize}
  \item identifier les capteurs utilisés dans la chaîne d'information d'un système ;
  \item identifier les constituants matériels de la chaîne d'information et les fonctions techniques réalisées ;
  \item définir la nature des informations d'entrée et de sortie d'un capteur ;
  \item justifier le choix d'un capteur au regard du cahier des charges ;
  \item extraire d'une documentation les valeurs numériques caractéristiques des solutions retenues ;
  \item analyser et mettre en œuvre une association de fonctions électroniques réalisant le conditionnement d'un signal ;
  \item décrire et représenter les évolutions temporelles et fréquentielles des signaux ;
  \item tracer le spectre d'un signal à partir de son expression temporelle et réaliser l'opération inverse ;
  \item déterminer la fonction de transfert d'un filtre passif ;
  \item prévoir le comportement d'un filtre aux basses et hautes fréquences ;
  \item tracer le spectre de sortie d'un filtre à partir du spectre d'entrée ;
  \item tracer le gabarit d'un filtre et déterminer l'ordre minimal associé ;
  \item déterminer la fréquence minimale d'échantillonnage ;
  \item déterminer le quantum, la résolution et le nombre de bits d'un CAN ou d'un CNA ;
  \item déterminer le gain d'un amplificateur permettant d'obtenir une résolution donnée ;
  \item coder une valeur en bases 2, 10 et 16 ainsi qu'en complément à deux.
\end{itemize}

\section{Chaîne d'acquisition}\label{sec:elec-chaine}

\subsection{Introduction}

Tous les systèmes pluritechniques ont besoin de contrôler de nombreux paramètres physiques : température, force, vitesse, position, déplacement, pression ou luminosité. À chacune de ces grandeurs peuvent correspondre un ou plusieurs types de capteurs, fondés sur des phénomènes physiques différents : variation de résistance, variation d'induction magnétique, effet Hall, effet piézoélectrique, etc.

Le capteur fournit un signal électrique — tension, courant ou variation d'impédance — image de la grandeur physique mesurée. Des circuits électroniques mettent ensuite en forme ce signal avant son exploitation par une unité de traitement : microcontrôleur, automate ou ordinateur.

\TLChaineAcquisition

La chaîne d'acquisition peut comporter les fonctions suivantes :
\begin{center}
\textbf{capteur et conditionneur $\rightarrow$ amplification $\rightarrow$ filtrage $\rightarrow$ CAN $\rightarrow$ unité de traitement $\rightarrow$ CNA.}
\end{center}

L'information fournie à l'unité de traitement peut se présenter sous trois formes :
\begin{itemize}
  \item signal analogique ;
  \item signal logique ;
  \item signal numérique.
\end{itemize}

\subsection{Nature du signal : analogique, logique et numérique}

Toutes les grandeurs qui nous entourent sont initialement analogiques : leur valeur est définie à chaque instant et peut prendre une infinité de valeurs. Cette forme est incompatible avec les seuls états « 0 » et « 1 » compris par une unité numérique.

Deux conversions sont alors possibles :
\begin{itemize}
  \item convertir l'information en un signal logique, ou signal tout ou rien ;
  \item convertir l'information en un signal numérique.
\end{itemize}

\subsubsection{Signal logique ou tout ou rien}

Un signal logique est obtenu en comparant le signal analogique à un seuil. Selon que la grandeur est inférieure ou supérieure à ce seuil, la sortie vaut « 0 » ou « 1 ». Cette transformation rend l'information directement exploitable par l'unité de traitement, mais elle fait perdre les détails relatifs aux variations de la grandeur analogique.

\paragraph{Exemple du destructeur d'aiguilles.}
Le destructeur d'aiguilles utilise un capteur infrarouge constitué d'une diode émettrice et d'une photodiode réceptrice. La présence ou l'absence de l'embase est interprétée par le microcontrôleur sous la forme d'un signal tout ou rien.

\subsubsection{Signal numérique}

Pour conserver davantage d'information qu'avec un signal logique, chaque valeur du signal analogique est codée par une suite de zéros et de uns appelée mot binaire.

Le signal est prélevé à intervalles réguliers : c'est l'échantillonnage. Cette opération entraîne une perte d'information, qui peut être limitée en choisissant une fréquence d'échantillonnage suffisante. Un compromis reste nécessaire : une fréquence trop élevée augmente la quantité de données à traiter et à mémoriser.

\subsection{Fonction « Traiter »}

La fonction « Traiter » dépend de la nature des informations :
\begin{itemize}
  \item les systèmes logiques traitent des informations bit à bit à l'aide de tables de vérité ou de graphes d'états ;
  \item les systèmes numériques traitent des mots binaires par programmation à partir d'algorithmes ou d'algorigrammes.
\end{itemize}

\section{Généralités sur les capteurs}\label{sec:elec-capteurs}

\subsection{Structure d'un capteur et vocabulaire utilisé}

Un capteur convertit une grandeur physique, appelée \emph{mesurande}, en un signal électrique image de ses variations. On parle aussi de transducteur. Dans la pratique, plusieurs conversions successives peuvent être nécessaires.

Le vocabulaire utilisé est le suivant :
\begin{description}
  \item[Mesurande :] grandeur physique que l'on souhaite mesurer : température, pression, force, position, etc.
  \item[Corps d'épreuve :] élément réagissant à la grandeur physique et fournissant une mesurande secondaire exploitable.
  \item[Transducteur :] élément lié au corps d'épreuve et transformant la mesurande secondaire en tension, courant ou variation d'impédance.
  \item[Conditionneur :] circuit électronique délivrant un signal exploitable à partir du signal brut du capteur. Il peut comporter une amplification, une mise en forme, un filtrage, un CAN ou un CNA.
\end{description}

\subsection{Différents capteurs}

Les systèmes mettent en œuvre de nombreux capteurs :
\begin{itemize}
  \item déplacements : potentiomètres résistifs, capteurs inductifs ou capacitifs, capteurs de proximité ;
  \item efforts et déformations : jauges d'extensiométrie, capteurs piézoélectriques ;
  \item courant : capteurs à effet Hall ;
  \item vitesse et position : génératrices tachymétriques, codeurs incrémentaux ou absolus, potentiomètres, résolveurs, tachymètres optiques ;
  \item pression, humidité et couple ;
  \item température : thermocouples, thermistances CTN et CTP, sondes Pt100 ;
  \item optique : photorésistances et fibres optiques.
\end{itemize}

\subsubsection{Classification par la nature du signal de sortie}

On distingue les capteurs analogiques, logiques et numériques. Un capteur analogique délivre un signal variant continûment. Un capteur intelligent peut intégrer toute la chaîne d'acquisition et communiquer directement avec un réseau ou un bus de terrain.

\subsubsection{Classification par le principe de transduction}

On distingue :
\begin{itemize}
  \item les capteurs passifs, pour lesquels la mesurande agit sur la résistance, l'inductance ou la capacité du transducteur ;
  \item les capteurs actifs, qui produisent directement une tension, un courant ou une charge électrique.
\end{itemize}

Quelques effets physiques classiques sont résumés ci-dessous.

\begin{center}
\begin{tabular}{lll}
\toprule
Mesurande & Effet utilisé & Grandeur électrique obtenue \\
\midrule
Température & thermoélectricité & tension \\
Flux optique & photoémission & courant \\
Température & pyroélectricité & charge \\
Force, pression, accélération & piézoélectricité & charge \\
Position & effet Hall & tension \\
Vitesse & induction électromagnétique & tension \\
\bottomrule
\end{tabular}
\end{center}

\subsection{Caractéristiques d'un capteur}

Le choix d'un capteur résulte d'un compromis entre plusieurs caractéristiques :
\begin{itemize}
  \item étendue de mesure ;
  \item résolution ;
  \item sensibilité ;
  \item précision ;
  \item rapidité ;
  \item coût ;
  \item bande passante.
\end{itemize}

\begin{description}
  \item[Étendue de mesure :] valeurs extrêmes accessibles à la grandeur mesurée.
  \item[Résolution :] plus petite variation de la grandeur d'entrée que le capteur peut déceler.
  \item[Sensibilité :] variation du signal de sortie rapportée à une variation de la grandeur mesurée ; elle correspond localement à la dérivée de la caractéristique entrée-sortie.
  \item[Précision :] aptitude du capteur à fournir une valeur proche de la valeur vraie.
  \item[Rapidité :] temps nécessaire pour que la sortie réagisse à une variation de l'entrée.
  \item[Bande passante :] intervalle de fréquences dans lequel le capteur fournit une mesure exploitable.
\end{description}

\begin{center}
\fbox{\parbox{0.85\linewidth}{\centering\textbf{Il faut bien distinguer la résolution, qui caractérise la plus petite variation détectable, et la précision, qui caractérise la proximité de la valeur mesurée avec la valeur vraie.}}}
\end{center}

\subsection{Conditionnement}

Les signaux issus des capteurs sont souvent faibles, parfois de l'ordre de quelques millivolts. Ils doivent être amplifiés avant d'être transmis à l'unité de traitement. Des convertisseurs peuvent aussi modifier la nature du signal : courant-tension, fréquence-tension, charge-tension, etc.

\subsection{Capteurs et conditionneurs classiques}

Quelques associations usuelles doivent être connues :
\begin{itemize}
  \item jauges de contraintes associées à un pont de Wheatstone ;
  \item accéléromètres utilisés pour mesurer une accélération ou détecter le sens d'un mouvement ;
  \item sondes de température Pt100, CTN ou CTP, parfois associées à des résistances de linéarisation ;
  \item capteurs de courant à effet Hall, utiles pour réaliser une boucle d'asservissement en courant.
\end{itemize}

\subsubsection{Application : accéléromètre ADXL335}

On utilise un accéléromètre ADXL335 pour mesurer une accélération maximale $\gamma_{nc,\max}=\SI{1.2}{\metre\per\second\squared}$. Le composant est alimenté sous $\SI{3}{\volt}$ et consomme $\SI{350}{\micro\ampere}$. Sa sensibilité est de $\SI{300}{\milli\volt\per g}$, sa plage de linéarité est de $\pm 3g$ et sa sortie vaut $\SI{1.5}{\volt}$ pour une accélération nulle.

\begin{enumerate}
  \item Déterminer la plage de tension correspondant au domaine de linéarité.
  \item Convertir ce domaine de linéarité en $\si{\metre\per\second\squared}$.
  \item Déterminer la tension mesurée pour une accélération de $\SI{1.2}{\metre\per\second\squared}$.
  \item Calculer la puissance consommée par le dispositif.
\end{enumerate}

\subsubsection{Application : mesure de température}

On souhaite mesurer la température d'une pièce entre $\SI{0}{\degreeCelsius}$ et $\SI{50}{\degreeCelsius}$ à l'aide d'un capteur donné pour une plage de $-\SI{50}{\degreeCelsius}$ à $\SI{100}{\degreeCelsius}$.

\begin{enumerate}
  \item Définir les caractéristiques nécessaires au choix de ce capteur.
  \item Préciser les critères permettant de comparer plusieurs solutions.
  \item À partir de la loi $R(T)$ fournie par la documentation du capteur, déterminer s'il s'agit d'une thermistance CTN ou CTP.
\end{enumerate}

\section{Outils mathématiques pour la chaîne d'information : électrocinétique}\label{sec:elec-electrocinetique}

\subsection{Généralités}

L'énergie électrique nécessaire au fonctionnement d'un circuit est fournie par des sources d'électricité appelées générateurs. Elle est disponible sous deux formes :
\begin{itemize}
  \item les sources continues, ou DC : piles, accumulateurs, panneaux solaires ;
  \item les sources alternatives, ou AC : centrales électriques, éoliennes, alternateurs.
\end{itemize}

L'énergie électrique est principalement produite sous forme alternative à partir de centrales qui transforment une énergie nucléaire, thermique, hydraulique, éolienne ou marémotrice. Son unité est le joule. En électricité, on utilise également le kilowattheure :
\[
\SI{1}{\kilo\watt\hour}=\SI{3.6}{\mega\joule}.
\]

\subsection{Électrocinétique}

L'électrocinétique étudie la circulation des charges électriques dans un circuit ou réseau électrique. Un circuit élémentaire est constitué de conducteurs métalliques et de composants électriques. Il comporte au minimum un générateur, un récepteur et les conducteurs qui les relient.

L'analogie hydraulique permet d'interpréter le fonctionnement : une pompe crée une différence de pression et met le fluide en mouvement ; de même, un générateur crée une différence de potentiel et provoque la circulation des charges.

\subsection{Charge et courant}

L'intensité électrique est définie comme le débit de charges traversant une section du conducteur :
\[
i(t)=\frac{\mathrm{d}q(t)}{\mathrm{d}t}.
\]

Le courant s'exprime en ampères et la charge en coulombs. La charge élémentaire de l'électron vaut $e=-1{,}6\times10^{-19}\,\si{\coulomb}$. Le sens conventionnel du courant est opposé au déplacement des électrons dans un conducteur métallique.

Un courant se mesure avec un ampèremètre placé en série dans la branche étudiée. Pour un régime continu, l'intensité est constante. Pour un régime variable, elle dépend du temps.

\subsection{Potentiel et différence de potentiel}

Chaque point d'un circuit possède un potentiel électrique $V$, défini par rapport à un point de référence choisi arbitrairement, appelé masse ou potentiel nul.

La tension entre deux points $A$ et $B$ est la différence de potentiel :
\[
U_{AB}=V_A-V_B.
\]

Une tension se mesure avec un voltmètre placé en dérivation entre les deux points. L'orientation de la flèche fixe l'ordre des potentiels dans la relation.

\subsection{Dipôles}

Un dipôle est un élément électrique comportant deux bornes. Son comportement est caractérisé par une relation entre la tension $u(t)$ à ses bornes et le courant $i(t)$ qui le traverse.

On distingue notamment :
\begin{itemize}
  \item les dipôles passifs : résistance, condensateur, inductance ;
  \item les dipôles actifs : sources de tension et sources de courant ;
  \item les dipôles linéaires et non linéaires ;
  \item les dipôles récepteurs et générateurs selon leur échange de puissance avec le reste du circuit.
\end{itemize}

\subsection{Sources de tension et de courant}

Une source idéale de tension impose une tension indépendante du courant débité. Une source idéale de courant impose un courant indépendant de la tension à ses bornes.

\begin{center}
\begin{circuitikz}
  \draw (0,0) to[V,l=$E$] (0,3);
  \draw (3,0) to[I,l=$I_0$] (3,3);
  \node[below] at (0,0) {source de tension};
  \node[below] at (3,0) {source de courant};
\end{circuitikz}
\end{center}

Une source réelle de tension est modélisée par une source idéale en série avec une résistance interne. Une source réelle de courant est modélisée par une source idéale en parallèle avec une résistance interne.

\subsection{Énergie et puissance}

La puissance électrique instantanée est définie par :
\[
p(t)=u(t)i(t).
\]

Elle s'exprime en watts. Elle traduit la quantité d'énergie fournie ou consommée par unité de temps :
\[
W(t_1,t_2)=\int_{t_1}^{t_2}p(t)\,\mathrm{d}t.
\]

En régime continu, lorsque $U$ et $I$ sont constants :
\[
P=UI,
\qquad
W=P\Delta t.
\]

\subsection{Point de fonctionnement}

Chaque point de la caractéristique tension-courant d'un dipôle représente un fonctionnement possible. Le point de fonctionnement réel d'un circuit correspond à l'intersection de la caractéristique du générateur et de celle du récepteur. Il fixe simultanément la tension aux bornes du récepteur et le courant qui le traverse.

\subsection{Conventions récepteur et générateur}

Le courant et la tension sont des grandeurs algébriques. Leur signe dépend de l'orientation choisie pour les flèches.

\begin{itemize}
  \item En convention récepteur, les flèches de tension et de courant sont de sens opposés.
  \item En convention générateur, les flèches de tension et de courant sont de même sens.
\end{itemize}

La convention est choisie en fonction du fonctionnement attendu. Une résistance fonctionne toujours comme récepteur. Un condensateur ou une inductance peuvent alternativement recevoir puis restituer de l'énergie.

\subsection{La résistance}

La résistance obéit à la loi d'Ohm :
\[
u(t)=Ri(t),
\qquad
R\text{ en }\si{\ohm}.
\]

La conductance vaut $G=1/R$ et s'exprime en siemens :
\[
i(t)=Gu(t).
\]

\subsubsection{Associations de résistances}

Pour des résistances en série :
\[
R_{\mathrm{eq}}=R_1+R_2+\cdots+R_n.
\]

Pour des résistances en parallèle :
\[
\frac{1}{R_{\mathrm{eq}}}=\frac{1}{R_1}+\frac{1}{R_2}+\cdots+\frac{1}{R_n}.
\]

Dans le cas de deux résistances en parallèle :
\[
R_{\mathrm{eq}}=\frac{R_1R_2}{R_1+R_2}.
\]

La puissance dissipée par effet Joule est :
\[
P_R=Ri^2=\frac{u^2}{R}=ui.
\]

\subsection{Le condensateur}

Un condensateur stocke de l'énergie sous forme électrostatique. Sa relation courant-tension est :
\[
i(t)=C\frac{\mathrm{d}u(t)}{\mathrm{d}t}.
\]

La tension aux bornes d'un condensateur ne peut pas varier instantanément. L'énergie emmagasinée vaut :
\[
W_C=\frac{1}{2}Cu^2.
\]

En régime continu établi, le condensateur se comporte comme un circuit ouvert.

Pour des condensateurs en parallèle :
\[
C_{\mathrm{eq}}=C_1+C_2+\cdots+C_n.
\]

Pour des condensateurs en série :
\[
\frac{1}{C_{\mathrm{eq}}}=\frac{1}{C_1}+\frac{1}{C_2}+\cdots+\frac{1}{C_n}.
\]

\subsection{L'inductance}

Une inductance stocke de l'énergie sous forme magnétique. Sa relation courant-tension est :
\[
u(t)=L\frac{\mathrm{d}i(t)}{\mathrm{d}t}.
\]

Le courant dans une inductance ne peut pas varier instantanément. L'énergie emmagasinée vaut :
\[
W_L=\frac{1}{2}Li^2.
\]

En régime continu établi, une inductance idéale se comporte comme un court-circuit.

Pour des inductances non couplées en série :
\[
L_{\mathrm{eq}}=L_1+L_2+\cdots+L_n.
\]

Pour des inductances non couplées en parallèle :
\[
\frac{1}{L_{\mathrm{eq}}}=\frac{1}{L_1}+\frac{1}{L_2}+\cdots+\frac{1}{L_n}.
\]

\subsection{Quadrants de fonctionnement}

Les conventions de fléchage ne déterminent pas à elles seules le mode de fonctionnement du dipôle. Celui-ci dépend du signe de la puissance $p=ui$.

En convention récepteur :
\begin{itemize}
  \item si $p>0$, le dipôle reçoit de l'énergie : il fonctionne en récepteur ;
  \item si $p<0$, le dipôle fournit de l'énergie : il fonctionne en générateur.
\end{itemize}

En convention générateur :
\begin{itemize}
  \item si $p>0$, le dipôle fonctionne en générateur ;
  \item si $p<0$, il fonctionne en récepteur.
\end{itemize}

Pour une résistance, la caractéristique $u=Ri$ passe par l'origine. En convention récepteur, les quadrants 1 et 3 correspondent au fonctionnement récepteur. Pour une source de tension idéale représentée en convention générateur, les quadrants sont interprétés de manière inverse.

\section{Outils et lois de l'électrocinétique}\label{sec:elec-lois}

\subsection{Lois de Kirchhoff}

Un circuit peut être constitué d'un ensemble de dipôles associés. On définit :
\begin{description}
  \item[Nœud :] connexion entre au moins deux dipôles ;
  \item[Branche :] portion de circuit comprise entre deux nœuds consécutifs ;
  \item[Maille :] chemin fermé constitué de plusieurs branches.
\end{description}

\subsubsection{Loi des nœuds}

La somme algébrique des courants arrivant à un nœud est nulle :
\[
\sum_k i_k=0.
\]

On peut aussi écrire que la somme des courants entrants est égale à la somme des courants sortants.

\subsubsection{Loi des mailles}

Dans une maille orientée, la somme algébrique des tensions est nulle :
\[
\sum_k u_k=0.
\]

Le signe d'une tension dépend de l'orientation de sa flèche par rapport au sens choisi pour parcourir la maille.

\subsection{Diviseur de tension}

Le diviseur de tension s'applique à des résistances placées en série et parcourues par le même courant.

\TLPontDiviseur

Pour deux résistances :
\[
u_s=E\frac{R_2}{R_1+R_2}.
\]

Plus généralement, la tension aux bornes de $R_k$ dans une association série vaut :
\[
U_k=U\frac{R_k}{\sum_i R_i}.
\]

\subsection{Diviseur de courant}

Le diviseur de courant s'applique à des résistances en parallèle soumises à la même tension.

Pour deux résistances parcourues par les courants $I_1$ et $I_2$, avec $I=I_1+I_2$ :
\[
I_1=I\frac{R_2}{R_1+R_2},
\qquad
I_2=I\frac{R_1}{R_1+R_2}.
\]

En utilisant les conductances :
\[
I_k=I\frac{G_k}{\sum_i G_i}.
\]

\subsection{Théorème de superposition}

Dans un circuit linéaire comportant plusieurs sources indépendantes, la tension ou le courant dans une branche est la somme algébrique des contributions obtenues en ne conservant qu'une source active à la fois.

Pour neutraliser les autres sources :
\begin{itemize}
  \item une source idéale de tension est remplacée par un court-circuit ;
  \item une source idéale de courant est remplacée par un circuit ouvert.
\end{itemize}

Le théorème s'applique aux tensions et aux courants, mais pas directement aux puissances.

\subsection{Théorème de Millman}

Le théorème de Millman permet de déterminer le potentiel d'un nœud relié à plusieurs branches contenant chacune une source de tension et une résistance.

Pour un nœud de potentiel $V$ relié à des potentiels $E_k$ par des résistances $R_k$ :
\[
V=\frac{\displaystyle\sum_k \frac{E_k}{R_k}}
        {\displaystyle\sum_k \frac{1}{R_k}}.
\]

Avec les conductances $G_k=1/R_k$ :
\[
V=\frac{\sum_k G_kE_k}{\sum_k G_k}.
\]

Le signe de chaque terme dépend de l'orientation de la source dans la branche correspondante.

\subsection{Théorème de Kennelly}

Le théorème de Kennelly permet de transformer un réseau de trois résistances montées en triangle en un réseau équivalent en étoile, et réciproquement.

Pour le passage triangle $R_{AB}$, $R_{BC}$, $R_{CA}$ vers étoile $R_A$, $R_B$, $R_C$ :
\[
R_A=\frac{R_{AB}R_{CA}}{R_{AB}+R_{BC}+R_{CA}},
\]
\[
R_B=\frac{R_{AB}R_{BC}}{R_{AB}+R_{BC}+R_{CA}},
\]
\[
R_C=\frac{R_{BC}R_{CA}}{R_{AB}+R_{BC}+R_{CA}}.
\]

Pour le passage étoile vers triangle :
\[
R_{AB}=R_A+R_B+\frac{R_AR_B}{R_C},
\]
\[
R_{BC}=R_B+R_C+\frac{R_BR_C}{R_A},
\]
\[
R_{CA}=R_C+R_A+\frac{R_CR_A}{R_B}.
\]

\subsection{Synthèse électrocinétique}

Pour choisir rapidement une méthode de résolution :
\begin{itemize}
  \item une source avec une association série : diviseur de tension ;
  \item une source avec une association parallèle : diviseur de courant ;
  \item plusieurs sources de tension reliées à un même nœud : théorème de Millman ;
  \item un mélange de sources de tension et de courant dans un circuit linéaire : superposition ;
  \item un réseau résistif non réductible par série-parallèle : transformation de Kennelly ou lois de Kirchhoff.
\end{itemize}

\subsubsection{Petites applications}

À partir des circuits proposés dans le document de cours :
\begin{enumerate}
  \item déterminer l'expression d'un potentiel de nœud $V_1$ ;
  \item déterminer l'expression d'un potentiel de nœud $V_4$ ;
  \item en déduire la différence de potentiel aux bornes d'une résistance donnée ;
  \item justifier le choix de la méthode utilisée parmi Kirchhoff, diviseur, Millman et superposition.
\end{enumerate}

% -----------------------------------------------------------------------------
% Les parties suivantes sont encore provisoires. Elles seront remplacées dans
% l'ordre du document source lors des prochaines reprises.
% -----------------------------------------------------------------------------

\section{Amplificateur opérationnel}\label{sec:elec-aop}

Un amplificateur opérationnel idéal présente un gain différentiel infini, des courants d'entrée nuls et une résistance de sortie nulle. En régime linéaire avec contre-réaction négative,
\[
i_+=i_-=0
\qquad\text{et}\qquad
u_+=u_-.
\]

\subsection{Montage inverseur}
\TLAOPInverseur
\[
\boxed{\frac{u_s}{u_e}=-\frac{R_2}{R_1}}.
\]

\subsection{Montage non-inverseur}
\TLAOPNonInverseur
\[
\boxed{\frac{u_s}{u_e}=1+\frac{R_2}{R_1}}.
\]

\section{Filtrage analogique}\label{sec:elec-filtrage}

\subsection{Filtre RC passe-bas}
\TLFiltreRCPasseBas
\[
H(p)=\frac{1}{1+RCp},
\qquad
f_c=\frac{1}{2\pi RC}.
\]

\subsection{Filtre RC passe-haut}
\TLFiltreRCPasseHaut
\[
H(p)=\frac{RCp}{1+RCp}.
\]

\section{Échantillonnage}\label{sec:elec-echantillonnage}

\TLEchantillonnageSignal
\[
f_e=\frac{1}{T_e},
\qquad
\boxed{f_e>2f_{\max}}.
\]

\section{Conversion analogique-numérique}\label{sec:elec-can}

Pour un convertisseur sur $N$ bits et une plage $[U_{\min},U_{\max}]$,
\[
q=\frac{U_{\max}-U_{\min}}{2^N}.
\]

\section{Conversion numérique-analogique}\label{sec:elec-cna}

Un convertisseur numérique-analogique produit une tension ou un courant proportionnel au mot binaire appliqué en entrée.

\section{À retenir}\label{sec:elec-retenir}

\begin{itemize}
  \item Une chaîne d'acquisition associe capteur, conditionnement, filtrage, conversion et traitement.
  \item La nature analogique, logique ou numérique de l'information conditionne son traitement.
  \item Le choix d'un capteur repose sur son étendue, sa résolution, sa précision, sa sensibilité, sa rapidité et sa bande passante.
  \item Le conditionnement adapte le signal du capteur à l'unité de traitement.
  \item Les lois de Kirchhoff, les diviseurs, Millman, la superposition et Kennelly constituent les principaux outils de résolution des circuits résistifs.
\end{itemize}


% Parties 5 à 9 : spectre, filtrage, numérisation, numération et sources.
% ============================================================================
% COURS 07 — PARTIES 5 À 10
% Reprise du document 11-AcquisitionTraitementSignal2022.docx
% ============================================================================

\section{Représentation spectrale}\label{sec:elec-spectrale}

\subsection{Intérêt de la représentation spectrale}

La représentation temporelle décrit l'évolution d'un signal en fonction du temps. La représentation spectrale décrit la répartition de son contenu en fréquence. Ces deux descriptions sont complémentaires : un même signal peut être observé soit dans le domaine temporel, soit dans le domaine fréquentiel.

La représentation spectrale permet notamment :
\begin{itemize}
  \item d'identifier la valeur moyenne d'un signal ;
  \item de déterminer sa fréquence fondamentale ;
  \item de repérer les harmoniques ;
  \item de connaître la bande fréquentielle occupée ;
  \item de prévoir l'effet d'un filtre sur le signal.
\end{itemize}

\subsection{Dualité temps-fréquence}

À un signal temporel $s(t)$ correspond une représentation fréquentielle $S(f)$. Une variation rapide dans le temps nécessite généralement des composantes de fréquence élevée. À l'inverse, un signal lentement variable possède principalement des composantes de basse fréquence.

Un signal sinusoïdal pur
\[
s(t)=A\cos(2\pi f_0t+\varphi)
\]
possède une seule fréquence $f_0$. Son spectre en amplitude comporte donc une raie à la fréquence $f_0$.

Une composante continue $S_0$ correspond à une raie à la fréquence nulle.

\subsection{Représentation spectrale d'un signal périodique}

Tout signal périodique de période $T$ et de fréquence fondamentale
\[
f_0=\frac{1}{T}
\]
peut, sous certaines conditions, être décomposé en une somme de sinusoïdes : c'est la décomposition en série de Fourier.

On peut écrire :
\[
s(t)=S_0+\sum_{n=1}^{+\infty}A_n\cos\left(2\pi n f_0t+\varphi_n\right).
\]

Dans cette expression :
\begin{itemize}
  \item $S_0$ est la valeur moyenne ;
  \item la composante de fréquence $f_0$ est le fondamental ;
  \item les composantes de fréquence $nf_0$ sont les harmoniques de rang $n$ ;
  \item $A_n$ est l'amplitude de l'harmonique de rang $n$ ;
  \item $\varphi_n$ est sa phase à l'origine.
\end{itemize}

Le spectre d'un signal périodique est un spectre discret : il est constitué de raies séparées de $f_0$.

\subsubsection{Signal sinusoïdal avec valeur moyenne}

Pour
\[
s(t)=S_0+A\cos(2\pi f_0t+\varphi),
\]
le spectre en amplitude contient :
\begin{itemize}
  \item une raie d'amplitude $S_0$ à $f=0$ ;
  \item une raie d'amplitude $A$ à $f=f_0$.
\end{itemize}

\subsubsection{Signal carré symétrique}

Un signal carré symétrique de valeur moyenne nulle contient uniquement des harmoniques impaires. Son développement comporte des termes aux fréquences
\[
f_0,\;3f_0,\;5f_0,\ldots
\]
avec des amplitudes décroissantes lorsque le rang augmente.

\subsubsection{Lecture d'un spectre}

À partir d'un spectre en amplitude, il faut savoir déterminer :
\begin{itemize}
  \item la valeur moyenne ;
  \item la fréquence fondamentale ;
  \item la période du signal ;
  \item le rang de chaque harmonique ;
  \item l'expression temporelle lorsque les phases sont connues ou nulles.
\end{itemize}

\begin{exemple}[Lecture d'un spectre]
Un spectre présente une raie de $\SI{2}{\volt}$ à $f=0$, une raie de $\SI{3}{\volt}$ à $\SI{50}{\hertz}$ et une raie de $\SI{1}{\volt}$ à $\SI{150}{\hertz}$. En supposant les phases nulles :
\[
s(t)=2+3\cos(2\pi 50t)+\cos(2\pi 150t).
\]
La fréquence fondamentale est $\SI{50}{\hertz}$ et la raie à $\SI{150}{\hertz}$ est l'harmonique de rang 3.
\end{exemple}

\subsection{Représentation spectrale d'un signal non périodique}

Un signal non périodique ne possède pas de fréquence fondamentale. Sa représentation fréquentielle est généralement continue : on parle de spectre continu.

Plus un signal est bref dans le temps, plus son spectre est étendu. Cette propriété explique qu'une impulsion très courte nécessite une bande passante importante pour être transmise sans déformation.

\begin{application}
Donner la représentation spectrale en amplitude des signaux temporels proposés dans le document d'origine. Pour chaque signal, préciser la valeur moyenne, le fondamental et le rang des harmoniques visibles.
\end{application}

\begin{application}
À partir des spectres en amplitude fournis, retrouver l'expression temporelle des signaux et indiquer leur valeur moyenne.
\end{application}

% ============================================================================
\section{Filtrage des signaux}\label{sec:elec-filtrage-complet}

\subsection{Définitions}

Un filtre est un système qui modifie le contenu fréquentiel d'un signal. Il laisse passer certaines fréquences et en atténue d'autres.

On distingue :
\begin{itemize}
  \item les filtres passifs, réalisés à partir de résistances, condensateurs et inductances ;
  \item les filtres actifs, qui utilisent en plus des composants actifs, notamment des amplificateurs opérationnels ;
  \item les filtres analogiques, qui agissent sur des signaux continus ;
  \item les filtres numériques, qui agissent sur des suites d'échantillons.
\end{itemize}

\subsection{Comparaison des types de filtre}

Les quatre types fondamentaux sont :
\begin{description}
  \item[Passe-bas :] transmet les basses fréquences et atténue les hautes fréquences.
  \item[Passe-haut :] atténue les basses fréquences et transmet les hautes fréquences.
  \item[Passe-bande :] transmet une bande de fréquences comprise entre deux fréquences de coupure.
  \item[Coupe-bande :] atténue une bande de fréquences et transmet les fréquences situées de part et d'autre.
\end{description}

\subsection{Gabarit d'un filtre}

Le gabarit définit les performances attendues du filtre. Il distingue :
\begin{itemize}
  \item la bande passante, dans laquelle le signal doit être peu atténué ;
  \item la bande atténuée, dans laquelle l'atténuation minimale est imposée ;
  \item la bande de transition, située entre les deux précédentes.
\end{itemize}

Pour un filtre passe-bas, on note généralement :
\begin{itemize}
  \item $f_p$ la limite de la bande passante ;
  \item $f_a$ le début de la bande atténuée ;
  \item $A_p$ l'atténuation maximale admise dans la bande passante ;
  \item $A_a$ l'atténuation minimale exigée dans la bande atténuée.
\end{itemize}

\subsection{Détermination de l'ordre à partir du gabarit}

L'ordre d'un filtre caractérise la rapidité avec laquelle son gain décroît dans la bande de transition. Pour un filtre passe-bas d'ordre $n$, la pente asymptotique est de
\[
-20n\ \text{dB par décade}.
\]

Une estimation de l'ordre minimal peut être obtenue à partir de la pente exigée entre deux points du gabarit :
\[
n\geq
\frac{A_a-A_p}{20\log_{10}(f_a/f_p)}.
\]

L'ordre retenu est l'entier immédiatement supérieur.

\subsection{Décibels}

Le gain en amplitude d'un filtre est
\[
G(f)=\left|\frac{S(f)}{E(f)}\right|.
\]

Le gain en décibels est défini par
\[
G_{\mathrm{dB}}(f)=20\log_{10}G(f).
\]

Quelques valeurs importantes :
\begin{center}
\begin{tabular}{cc}
\toprule
Gain en amplitude & Gain en décibels \\
\midrule
$1$ & $0$ dB \\
$\sqrt{2}$ & $+3$ dB \\
$1/\sqrt{2}$ & $-3$ dB \\
$10$ & $+20$ dB \\
$0{,}1$ & $-20$ dB \\
\bottomrule
\end{tabular}
\end{center}

Pour une puissance, le rapport en décibels s'écrit
\[
G_{\mathrm{dB}}=10\log_{10}\left(\frac{P_s}{P_e}\right).
\]

\subsection{Bande passante et fréquences de coupure}

La fréquence de coupure est généralement définie par une diminution de $\SI{3}{\decibel}$ par rapport au gain maximal. Elle correspond à un gain en amplitude divisé par $\sqrt{2}$.

La bande passante à $-3$ dB est l'ensemble des fréquences pour lesquelles
\[
G_{\mathrm{dB}}\geq G_{\mathrm{dB,max}}-3.
\]

Pour un passe-bande, la largeur de bande vaut
\[
\Delta f=f_{c2}-f_{c1}.
\]

\subsection{Quadripôles}

Un filtre peut être représenté comme un quadripôle possédant :
\begin{itemize}
  \item deux bornes d'entrée ;
  \item deux bornes de sortie ;
  \item une tension d'entrée $u_e(t)$ ;
  \item une tension de sortie $u_s(t)$.
\end{itemize}

Le comportement fréquentiel est décrit par la fonction de transfert complexe
\[
H(j\omega)=\frac{\underline{U_s}}{\underline{U_e}}.
\]

\subsection{Représentation complexe d'un signal sinusoïdal}

À la grandeur sinusoïdale
\[
u(t)=U_m\cos(\omega t+\varphi)
\]
on associe le nombre complexe
\[
\underline{U}=U_m e^{j\varphi}.
\]

La dérivation temporelle correspond à une multiplication par $j\omega$ :
\[
\frac{\mathrm d}{\mathrm dt}\longleftrightarrow j\omega.
\]

L'intégration correspond à une division par $j\omega$.

\subsection{Impédance et admittance complexes}

En régime sinusoïdal, l'impédance complexe d'un dipôle est définie par
\[
\underline{Z}=\frac{\underline{U}}{\underline{I}}.
\]

On écrit aussi
\[
\underline{Z}=R+jX,
\]
où $R$ est la résistance et $X$ la réactance.

L'admittance complexe est
\[
\underline{Y}=\frac{1}{\underline{Z}}=G+jB,
\]
où $G$ est la conductance et $B$ la susceptance.

Les impédances des dipôles élémentaires sont :
\[
\underline{Z_R}=R,
\qquad
\underline{Z_C}=\frac{1}{j\omega C},
\qquad
\underline{Z_L}=j\omega L.
\]

Les lois d'association restent valables avec les impédances complexes :
\begin{itemize}
  \item en série, les impédances s'ajoutent ;
  \item en parallèle, les admittances s'ajoutent.
\end{itemize}

\subsection{Fonction de transfert d'un filtre}

La fonction de transfert est
\[
H(j\omega)=\frac{\underline{U_s}}{\underline{U_e}}.
\]

Son module donne le gain en amplitude :
\[
G(\omega)=|H(j\omega)|.
\]

Son argument donne le déphasage de la sortie par rapport à l'entrée :
\[
\varphi(\omega)=\arg H(j\omega).
\]

\subsection{Méthode de détermination}

Pour déterminer la fonction de transfert d'un filtre passif :
\begin{enumerate}
  \item remplacer chaque dipôle par son impédance complexe ;
  \item repérer les associations série ou parallèle ;
  \item appliquer un diviseur de tension ou les lois de Kirchhoff ;
  \item former le rapport $\underline{U_s}/\underline{U_e}$ ;
  \item simplifier et mettre le résultat sous une forme canonique ;
  \item identifier le gain statique, la pulsation de coupure et l'ordre.
\end{enumerate}

\subsubsection{Filtre RC passe-bas}

\TLFiltreRCPasseBas

\[
H(j\omega)=\frac{1}{1+j\omega RC}
=\frac{1}{1+j\omega/\omega_c},
\qquad
\omega_c=\frac{1}{RC}.
\]

Le module et la phase sont :
\[
|H(j\omega)|=\frac{1}{\sqrt{1+(\omega RC)^2}},
\]
\[
\arg H(j\omega)=-\arctan(\omega RC).
\]

\subsubsection{Filtre RC passe-haut}

\TLFiltreRCPasseHaut

\[
H(j\omega)=\frac{j\omega RC}{1+j\omega RC}.
\]

\subsection{Prévision du comportement aux basses et hautes fréquences}

Une méthode rapide consiste à remplacer les condensateurs et inductances par leurs comportements limites.

Pour un condensateur :
\begin{itemize}
  \item à basse fréquence, $|Z_C|\to +\infty$ : il se comporte comme un circuit ouvert ;
  \item à haute fréquence, $|Z_C|\to 0$ : il se comporte comme un court-circuit.
\end{itemize}

Pour une inductance :
\begin{itemize}
  \item à basse fréquence, $|Z_L|\to 0$ : elle se comporte comme un court-circuit ;
  \item à haute fréquence, $|Z_L|\to +\infty$ : elle se comporte comme un circuit ouvert.
\end{itemize}

Cette étude permet d'identifier rapidement le type de filtre avant tout calcul détaillé.

\subsection{Mise en cascade de cellules élémentaires}

Lorsque plusieurs cellules sont mises en cascade sans interaction entre elles, la fonction de transfert globale est le produit des fonctions de transfert :
\[
H(j\omega)=H_1(j\omega)H_2(j\omega)\cdots H_n(j\omega).
\]

En décibels, les gains s'ajoutent :
\[
G_{\mathrm{dB}}=G_{1,\mathrm{dB}}+G_{2,\mathrm{dB}}+\cdots+G_{n,\mathrm{dB}}.
\]

Les phases s'ajoutent également.

\begin{application}
Donner les fonctions de transfert des filtres proposés dans le document d'origine, puis identifier leur type par l'étude des comportements aux basses et hautes fréquences.
\end{application}

% ============================================================================
\section{Échantillonnage et numérisation}\label{sec:elec-numerisation}

\subsection{Convertisseur analogique-numérique}

Un CAN transforme une tension analogique en un mot binaire. Pour un convertisseur sur $n$ bits, le nombre de codes disponibles est
\[
N=2^n.
\]

Pour une plage de conversion comprise entre $U_{\min}$ et $U_{\max}$, le quantum vaut
\[
q=\frac{U_{\max}-U_{\min}}{2^n}.
\]

Selon la convention retenue dans la documentation, le dénominateur peut être $2^n-1$ lorsque les deux valeurs extrêmes sont représentées exactement. Il faut toujours vérifier la définition donnée par le constructeur.

\subsection{Coder une information sur $n$ bits}

Le code numérique associé à une tension $u_e$ peut être déterminé par
\[
N=\left\lfloor\frac{u_e-U_{\min}}{q}\right\rfloor.
\]

Le résultat doit être limité à l'intervalle des codes admissibles. Une tension hors plage entraîne une saturation.

L'erreur de quantification est liée au remplacement d'une valeur continue par un niveau discret. Elle est généralement comprise entre
\[
-\frac{q}{2}\quad\text{et}\quad +\frac{q}{2}.
\]

\subsection{Technologies des CAN et choix}

Les technologies courantes sont :
\begin{itemize}
  \item CAN flash : très rapide mais nécessitant beaucoup de comparateurs ;
  \item CAN à approximations successives : bon compromis entre vitesse, résolution et coût ;
  \item CAN double rampe : lent mais précis, utilisé dans les appareils de mesure ;
  \item CAN sigma-delta : grande résolution pour des bandes passantes limitées.
\end{itemize}

Le choix dépend notamment :
\begin{itemize}
  \item du nombre de bits ;
  \item de la plage d'entrée ;
  \item du temps de conversion ;
  \item de la fréquence maximale d'échantillonnage ;
  \item de la précision ;
  \item du coût et de la consommation.
\end{itemize}

\subsection{Convertisseur numérique-analogique}

Un CNA transforme un mot binaire en une tension ou un courant analogique. Le quantum et la résolution se déterminent de manière analogue au CAN.

Après conversion, la sortie présente souvent une allure en escalier. Un filtre de reconstruction peut être utilisé pour lisser le signal.

\subsection{Échantillonnage}

L'échantillonnage consiste à prélever la valeur d'un signal analogique à intervalles réguliers $T_e$.

\[
f_e=\frac{1}{T_e}.
\]

\TLEchantillonnageSignal

Une fréquence d'échantillonnage élevée améliore la fidélité de la représentation, mais augmente la quantité de données et impose un temps de conversion plus court.

\subsection{Échantillonneur-bloqueur}

L'échantillonneur-bloqueur prélève la tension d'entrée puis la maintient constante pendant la durée nécessaire à la conversion. Il évite que la tension d'entrée varie pendant le fonctionnement du CAN.

\subsection{Théorème de Shannon-Nyquist}

Pour échantillonner sans ambiguïté un signal dont la fréquence maximale est $f_{\max}$, il faut respecter :
\[
\boxed{f_e>2f_{\max}}.
\]

Dans la pratique, on choisit souvent une fréquence sensiblement supérieure à cette limite.

\subsection{Repliement de spectre}

L'échantillonnage crée des répétitions du spectre autour des multiples de $f_e$. Si $f_e$ est insuffisante, ces répétitions se recouvrent : c'est le repliement de spectre. Des composantes hautes fréquences sont alors interprétées à tort comme des composantes de plus basse fréquence.

\subsection{Filtre anti-repliement}

Le filtre anti-repliement est un filtre passe-bas placé avant le CAN. Il élimine les composantes susceptibles de provoquer un recouvrement spectral.

Son gabarit doit être défini conjointement avec la fréquence d'échantillonnage. À partir du gabarit, on identifie la fréquence maximale utile $f_{\max}$ puis on impose
\[
f_e>2f_{\max}.
\]

\subsection{Déterminer la fréquence d'échantillonnage}

Le choix de $f_e$ doit satisfaire deux contraintes :
\begin{itemize}
  \item la condition de Shannon-Nyquist ;
  \item le temps de conversion du CAN.
\end{itemize}

Si le temps de conversion est $t_c$, il faut disposer d'une période d'échantillonnage compatible :
\[
T_e\geq t_c,
\qquad
f_e\leq \frac{1}{t_c}.
\]

Pour un signal sinusoïdal pur de fréquence $f$, on prend $f_{\max}=f$.

Pour un signal périodique non sinusoïdal, il faut considérer les harmoniques que l'on souhaite conserver. La fréquence maximale utile dépend alors du rang maximal retenu.

% ============================================================================
\section{Codage numérique des informations : numération}\label{sec:elec-numeration}

\subsection{Systèmes de numération}

Dans une base $b$, un nombre s'écrit à l'aide de chiffres compris entre $0$ et $b-1$. Sa valeur décimale est obtenue par
\[
(a_na_{n-1}\ldots a_0)_b=\sum_{k=0}^{n}a_kb^k.
\]

\subsection{Base binaire}

La base 2 utilise les chiffres 0 et 1. Chaque chiffre binaire est appelé bit.

Le bit de poids faible est le LSB et le bit de poids fort le MSB.

Par exemple :
\[
(101101)_2=1\times2^5+0\times2^4+1\times2^3+1\times2^2+0\times2+1=45.
\]

\subsection{Base hexadécimale}

La base 16 utilise les symboles
\[
0,1,\ldots,9,A,B,C,D,E,F.
\]

Chaque chiffre hexadécimal correspond exactement à quatre bits :
\[
(A)_{16}=(1010)_2,
\qquad
(F)_{16}=(1111)_2.
\]

\subsection{Notations}

Les notations courantes sont :
\begin{itemize}
  \item suffixe $_2$ pour le binaire ;
  \item suffixe $_{10}$ pour le décimal ;
  \item suffixe $_{16}$ ou préfixe \,\texttt{0x} pour l'hexadécimal.
\end{itemize}

\subsection{Changement de base}

Pour passer du binaire au décimal, on utilise les puissances de 2.

Pour passer du décimal au binaire, on effectue des divisions successives par 2 et on lit les restes dans l'ordre inverse.

Pour passer du binaire à l'hexadécimal, on regroupe les bits par paquets de quatre à partir de la droite.

\subsection{Addition en base 2}

Les règles élémentaires sont :
\[
0+0=0,
\quad
0+1=1,
\quad
1+1=10,
\quad
1+1+1=11.
\]

Une retenue dépassant le nombre de bits disponibles constitue un débordement.

\subsection{Nombres signés et complément à deux}

Sur $n$ bits, un nombre signé en complément à deux appartient à l'intervalle
\[
-2^{n-1}\leq N\leq 2^{n-1}-1.
\]

Pour coder l'opposé d'un nombre positif :
\begin{enumerate}
  \item écrire sa valeur binaire sur $n$ bits ;
  \item inverser tous les bits ;
  \item ajouter 1.
\end{enumerate}

Le bit de poids fort indique le signe : 0 pour un nombre positif, 1 pour un nombre négatif.

\subsection{Multiplication et division par $2^n$}

Un décalage de $n$ bits vers la gauche correspond, en l'absence de débordement, à une multiplication par $2^n$.

Un décalage de $n$ bits vers la droite correspond à une division entière par $2^n$. Pour les nombres signés, un décalage arithmétique conserve le bit de signe.

\subsection{Autres codes}

D'autres représentations peuvent être utilisées :
\begin{itemize}
  \item code BCD, qui code séparément chaque chiffre décimal ;
  \item code Gray, dans lequel deux valeurs successives ne diffèrent que d'un bit ;
  \item code ASCII pour représenter des caractères.
\end{itemize}

% ============================================================================
\section{Sources}\label{sec:elec-sources}

Cette partie reprend les références et documentations techniques utilisées dans le document d'origine : documentations de capteurs, convertisseurs, filtres et composants électroniques. Les références détaillées seront conservées lors de l'intégration des figures et tableaux d'origine.

% ============================================================================
\section{Exercices du chapitre}\label{sec:elec-exercices}

Les exercices du document source occupent les pages 93 à 135. Ils portent notamment sur :
\begin{itemize}
  \item l'analyse de chaînes d'acquisition ;
  \item le choix et le conditionnement de capteurs ;
  \item l'électrocinétique et les lois de Kirchhoff ;
  \item la représentation spectrale ;
  \item le calcul de fonctions de transfert de filtres ;
  \item les gabarits et l'ordre des filtres ;
  \item les CAN, CNA et l'échantillonnage ;
  \item la numération et le complément à deux.
\end{itemize}

\begin{application}
Les énoncés complets et leurs figures seront intégrés exercice par exercice à partir des pages 93 à 135 du document Word, afin de préserver les schémas et données numériques d'origine.
\end{application}

\section{Synthèse générale}\label{sec:elec-synthese-generale}

\begin{itemize}
  \item Un signal peut être étudié dans les domaines temporel et fréquentiel.
  \item Un signal périodique possède un spectre discret constitué du fondamental et de ses harmoniques.
  \item Un filtre est caractérisé par son gabarit, sa fonction de transfert, son gain et sa phase.
  \item Le comportement limite des condensateurs et inductances permet d'identifier rapidement un filtre.
  \item Le CAN quantifie et code les échantillons d'un signal analogique.
  \item La fréquence d'échantillonnage doit respecter Shannon-Nyquist et le temps de conversion.
  \item Le filtre anti-repliement limite le spectre avant le CAN.
  \item Les bases 2 et 16 et le complément à deux sont indispensables au traitement numérique.
\end{itemize}


% Partie 10 : exercices du chapitre réécrits à partir du document source.
% ============================================================================
% COURS 07 — EXERCICES DU CHAPITRE
% Reprise progressive des exercices du document 11-AcquisitionTraitementSignal2022.docx
% ============================================================================

\section{Exercices du chapitre}\label{sec:elec-exercices}

\subsection{Capteurs classiques}

Pour chacun des capteurs suivants, donner son rôle, préciser la grandeur physique mesurée, indiquer la nature du signal de sortie et expliquer son principe de fonctionnement à l'aide d'un schéma de principe :
\begin{enumerate}
  \item génératrice tachymétrique ;
  \item potentiomètre rotatif ;
  \item codeur incrémental ;
  \item codeur absolu.
\end{enumerate}

Pour chaque solution, préciser également :
\begin{itemize}
  \item s'il s'agit d'un capteur actif ou passif ;
  \item si le signal délivré est analogique, logique ou numérique ;
  \item les principales caractéristiques à vérifier pour choisir le capteur : étendue de mesure, résolution, précision, rapidité et bande passante.
\end{itemize}

\subsection{Le banc Balafre — chaîne d'acquisition d'efforts}

Entre autres contrôles de la chaîne d'acquisition, le superviseur vérifie que la mesure des efforts se fait correctement au niveau des actionneurs piézoélectriques et au niveau du joint testé. Les capteurs de force utilisés sur le système sont analogiques.

Afin de simplifier le traitement et l'interprétation de ces forces, on utilise un amplificateur de charge à plusieurs canaux. Cet amplificateur possède notamment :
\begin{itemize}
  \item un amplificateur de sommation pour le calcul analogique des forces et moments résultants ;
  \item un convertisseur analogique-numérique pour le traitement des données par l'algorithme de contrôle.
\end{itemize}

Dans l'algorithme de contrôle, la valeur d'effort de chaque actionneur est comparée à la valeur théorique de la consigne. Si un écart trop important est constaté, l'algorithme émet un signal d'erreur.

Pour cette mesure, on considère qu'une résolution inférieure à \SI{10}{\newton} est nécessaire. La conversion analogique-numérique est réalisée sur 12 bits. La mesure de l'effort s'effectue sur la plage de $-\SI{20}{\kilo\newton}$ à $+\SI{20}{\kilo\newton}$.

\begin{enumerate}
  \item Représenter la chaîne d'acquisition depuis le capteur piézoélectrique jusqu'à l'unité de traitement.
  \item Identifier les fonctions assurées par l'amplificateur de charge.
  \item Calculer le nombre de niveaux de quantification du CAN.
  \item Déterminer le quantum correspondant à la plage complète de mesure.
  \item Vérifier que la résolution obtenue respecte l'exigence de \SI{10}{\newton}.
  \item Déterminer le nombre minimal de bits nécessaire pour respecter cette exigence sur la même plage de mesure.
  \item Expliquer pourquoi une amplification et un conditionnement sont nécessaires avant la conversion analogique-numérique.
\end{enumerate}

\subsection{Accéléromètre ADXL335}

On utilise un accéléromètre ADXL335 pour mesurer une accélération maximale $\gamma_{\max}=\SI{1.2}{\metre\per\second\squared}$. Le composant est alimenté sous \SI{3}{\volt} et consomme \SI{350}{\micro\ampere}. Sa sensibilité est de \SI{300}{\milli\volt\per g}, sa plage de linéarité est de $\pm 3g$ et sa tension de sortie vaut \SI{1.5}{\volt} pour une accélération nulle.

\begin{enumerate}
  \item Déterminer la plage de tension de sortie correspondant au domaine de linéarité.
  \item Convertir le domaine de linéarité en \si{\metre\per\second\squared}.
  \item Déterminer la tension de sortie pour une accélération de \SI{1.2}{\metre\per\second\squared}.
  \item Calculer la puissance électrique consommée par le dispositif.
  \item Proposer une chaîne de conditionnement permettant d'utiliser ce capteur avec un CAN dont la plage d'entrée est $[0;\SI{3.3}{\volt}]$.
\end{enumerate}

\subsection{Mesure de température}

On souhaite mesurer la température d'une pièce entre \SI{0}{\degreeCelsius} et \SI{50}{\degreeCelsius}. Le capteur disponible possède une plage de mesure allant de $-\SI{50}{\degreeCelsius}$ à \SI{100}{\degreeCelsius}.

\begin{enumerate}
  \item Définir l'étendue de mesure du capteur et celle réellement utile pour l'application.
  \item Préciser les critères permettant de comparer plusieurs capteurs possibles.
  \item Expliquer la différence entre précision et résolution dans cette application.
  \item À partir d'une caractéristique $R(T)$ décroissante, déterminer s'il s'agit d'une thermistance CTN ou CTP.
  \item Proposer un montage de conditionnement simple permettant de transformer la variation de résistance en tension.
  \item Indiquer les précautions à prendre pour limiter l'auto-échauffement du capteur.
\end{enumerate}

\subsection{Numérisation d'une grandeur analogique}

Un capteur fournit une tension comprise entre \SI{0.5}{\volt} et \SI{4.5}{\volt}. Cette tension est convertie par un CAN sur 10 bits dont la plage d'entrée est $[0;\SI{5}{\volt}]$.

\begin{enumerate}
  \item Calculer le nombre de niveaux du CAN.
  \item Déterminer le quantum du convertisseur.
  \item Déterminer les codes numériques correspondant aux tensions \SI{0.5}{\volt}, \SI{2.5}{\volt} et \SI{4.5}{\volt}.
  \item Calculer la résolution réellement obtenue sur la plage utile du capteur.
  \item Proposer une adaptation affine permettant d'utiliser toute la plage $[0;\SI{5}{\volt}]$ du CAN.
  \item Déterminer le gain et le décalage nécessaires pour cette adaptation.
\end{enumerate}

\subsection{Choix d'une fréquence d'échantillonnage}

Un signal utile contient des composantes fréquentielles jusqu'à \SI{800}{\hertz}. Il est acquis par une chaîne comprenant un filtre anti-repliement puis un CAN.

\begin{enumerate}
  \item Déterminer la fréquence minimale d'échantillonnage imposée par le théorème de Shannon-Nyquist.
  \item Proposer une fréquence d'échantillonnage pratique et justifier le choix.
  \item Expliquer le rôle du filtre anti-repliement.
  \item Indiquer ce qui se produit si une composante de fréquence \SI{1.4}{\kilo\hertz} est échantillonnée à \SI{2}{\kilo\hertz} sans filtrage préalable.
  \item Préciser la relation entre période d'échantillonnage et temps de conversion maximal du CAN.
\end{enumerate}

\subsection{Filtrage d'un signal de capteur}

Un capteur délivre un signal utile de fréquence inférieure à \SI{50}{\hertz}, auquel se superpose un bruit important à \SI{1}{\kilo\hertz}. On utilise un filtre RC passe-bas.

\begin{enumerate}
  \item Justifier le choix d'un filtre passe-bas.
  \item Établir la fonction de transfert du montage.
  \item Choisir une fréquence de coupure compatible avec les données du problème.
  \item Pour $R=\SI{10}{\kilo\ohm}$, déterminer la valeur de $C$ correspondante.
  \item Calculer l'atténuation du bruit à \SI{1}{\kilo\hertz}.
  \item Vérifier que le signal utile à \SI{50}{\hertz} reste suffisamment peu atténué.
\end{enumerate}

\subsection{Codeur incrémental}

Un codeur incrémental optique est utilisé pour mesurer la position angulaire d'un arbre.

\begin{enumerate}
  \item Donner le rôle et le principe de fonctionnement d'un codeur incrémental optique. Illustrer la réponse par un schéma de principe faisant apparaître le disque, la source lumineuse et le photodétecteur.
  \item Le codeur est équipé d'une seule voie de mesure et d'un disque comportant 25 fentes. Déterminer la résolution angulaire du capteur en degrés.
  \item Déterminer la nouvelle résolution lorsque le codeur est équipé de deux voies de mesure en quadrature.
\end{enumerate}

Le codeur est monté sur l'arbre d'un moteur suivi d'un réducteur de rapport 100.

\begin{enumerate}
  \setcounter{enumi}{3}
  \item Déterminer la résolution obtenue vis-à-vis de l'arbre de sortie du réducteur.
\end{enumerate}

La position mesurée est ensuite transformée par un convertisseur numérique-analogique. Ce convertisseur associe une plage angulaire de $-10$ tours à $+10$ tours à une plage de tension de $-\SI{5}{\volt}$ à $+\SI{5}{\volt}$.

\begin{enumerate}
  \setcounter{enumi}{4}
  \item Déterminer le gain du convertisseur numérique-analogique, exprimé en volts par tour puis en volts par radian.
\end{enumerate}

% ============================================================================
% COURS 07 — EXERCICES COMPLÉMENTAIRES
% Reprise du document 11-AcquisitionTraitementSignal2022.docx
% ============================================================================

\subsection{Pompe médicale OPTIMA3 — acquisition d'une pression}

La pompe médicale OPTIMA3 est une pompe volumétrique utilisée en milieu hospitalier. Elle assure l'injection d'un liquide par perfusion sans intervention permanente du personnel médical.

Une anomalie possible est l'occlusion de la tubulure : le canal peut être fermé par écrasement, pliage ou obstruction de la veine. Dans ce cas, la pompe continue à débiter et la pression augmente dangereusement. Un capteur de pression est donc placé au contact de la tubulure. Au-delà d'un seuil étalonné en usine, le microcontrôleur arrête le moteur.

La partie sensible du capteur est réalisée dans un matériau souple dans lequel est noyé un fil électrique. La déformation de ce fil modifie sa résistivité et donc la tension délivrée par le capteur.

\begin{enumerate}
  \item Identifier la nature de l'information délivrée par le capteur.
  \item À partir de la caractéristique entrée-sortie fournie dans le document technique, calculer la sensibilité du capteur en \si{\milli\volt\per\mega\pascal}.
\end{enumerate}

L'information est envoyée au microcontrôleur qui possède un convertisseur analogique-numérique. La résolution souhaitée sur la pression doit rester inférieure à \SI{0.01}{\mega\pascal}.

La chaîne d'acquisition comprend le capteur, un conditionneur, un filtre anti-repliement, le CAN et l'unité de traitement.

\begin{enumerate}
  \setcounter{enumi}{2}
  \item Calculer le gain de l'amplificateur du conditionneur afin d'obtenir une augmentation de \SI{1}{\volt} en sortie pour une augmentation de pression de \SI{0.1}{\mega\pascal}.
  \item Déterminer la plage de tension en sortie du conditionneur pour l'étendue de pression utile.
  \item À partir du nombre de bits et de la pleine échelle du CAN donnés dans le document technique, déterminer son quantum en volts.
  \item En déduire la résolution obtenue sur la mesure de pression et vérifier l'exigence de \SI{0.01}{\mega\pascal}.
  \item À partir du spectre du signal de pression, déterminer la fréquence maximale utile.
  \item En déduire la fréquence minimale d'échantillonnage imposée par Shannon-Nyquist.
  \item Proposer le gabarit du filtre anti-repliement placé avant le CAN.
\end{enumerate}

\subsection{Mesure du courant image du couple moteur}

Le couple d'un moteur est estimé à partir de la mesure de son courant par un capteur à effet Hall. Le capteur délivre une tension image du courant moteur. Cette tension comporte une composante moyenne utile et une ondulation liée à la commande du moteur.

On souhaite conserver l'image du courant moyen et atténuer suffisamment l'ondulation afin de disposer d'une mesure exploitable pour l'asservissement. Un filtre passif est utilisé et son étude est menée avec la représentation complexe.

\begin{enumerate}
  \item Indiquer le type de filtre permettant de conserver la composante continue tout en atténuant l'ondulation.
  \item Pour la structure RC proposée dans le document, exprimer la tension de sortie en fonction de la tension d'entrée et des impédances des composants.
  \item Donner l'impédance complexe du condensateur.
  \item En déduire la fonction de transfert du filtre et la mettre sous la forme canonique
  \[
    H(j\omega)=\frac{1}{1+j\dfrac{\omega}{\omega_c}},
    \qquad \omega_c=\frac{1}{RC}.
  \]
  \item À partir de l'atténuation demandée au fondamental de l'ondulation et de sa fréquence, déterminer la pulsation de coupure maximale admissible.
  \item Afin de ne pas charger excessivement le capteur à effet Hall, l'impédance d'entrée du montage ne doit pas être inférieure à la valeur imposée dans le document. En déduire la valeur minimale de la résistance $R$.
  \item Calculer ensuite la capacité $C$ permettant de respecter la pulsation de coupure déterminée précédemment.
  \item Vérifier le gain du filtre à la fréquence de l'ondulation et conclure sur le respect de l'atténuation demandée.
\end{enumerate}

\subsection{Chaîne d'acquisition d'un effort sur une barre}

En fonctionnement normal, l'effort exercé sur une barre doit être compris entre \SI{5}{\newton} et \SI{100}{\newton}. Un effort inférieur signale une mauvaise mise en place ou un actionneur défectueux ; un effort supérieur peut provenir d'un ralentissement trop brutal ou d'une action volontaire des passagers.

Le capteur est constitué d'un pont de jauges collé sur la barre. Sa sensibilité est de \SI{1.5}{\milli\volt\per\volt}, son étendue de mesure est de \SI{250}{\newton} et il est alimenté sous \SI{5}{\volt}. Le signal délivré vaut alors \SI{30}{\micro\volt} par newton.

Le CAN de l'automate possède les caractéristiques suivantes :
\begin{itemize}
  \item quantum en tension : \SI{2.5}{\milli\volt} ;
  \item codage sur 12 bits ;
  \item fréquence d'échantillonnage : \SI{500}{\hertz}.
\end{itemize}

\begin{enumerate}
  \item Justifier la zone d'implantation des jauges sur la barre.
  \item Déterminer la tension brute délivrée par le pont de jauges pour \SI{5}{\newton}, \SI{100}{\newton} et \SI{250}{\newton}.
  \item Montrer que l'entrée analogique de l'automate ne peut pas exploiter directement ce signal.
  \item Déterminer le gain minimal du conditionneur permettant d'obtenir une résolution inférieure ou égale à \SI{1}{\newton}.
  \item Vérifier l'absence de saturation pour l'effort maximal de \SI{250}{\newton} et la pleine échelle du CAN.
  \item Déterminer la période d'échantillonnage du convertisseur.
  \item Donner la fréquence maximale théorique du signal mesurable conformément au théorème de Shannon-Nyquist.
\end{enumerate}


\printindex
\end{document}
